\section{Inertia and descent over a torus}\label{sec:descent}

We attach boundary coefficients to the orders of meridians in a fixed
group quotient. These orders also control components whose images are
exceptional for an intermediate map. Keeping them on every model will
allow us to descend positivity from the fibres over a torus.

\subsection{Meridian orders and differential pullback}

Let \(c:M\dashrightarrow Z\) be a fibration from a smooth compact
K\"ahler manifold, and let
\(\rho:\pi_1(M)\twoheadrightarrow R\) kill the image of the group
of a smooth general \(c\)-fibre. Replace the source and target by
smooth models so that \(c\) is holomorphic, and choose an analytic
Zariski open \(Z^\circ\) over which it is smooth. The homotopy
sequence of the smooth fibration gives a homomorphism
\[
 \psi:\pi_1(Z^\circ)\longrightarrow R
\]
compatible with \(\rho\). It is surjective: the open source maps
surjectively on the fundamental group of the compact source. The
homomorphism is independent of a further shrinking, in this
compatibility sense.

For a prime divisor \(D\) on the smooth base, let \(m_D\) be the
order in \(R\) of a meridian about its general point. Meridians are
defined up to conjugacy, which does not affect their orders. Set
\(m_D=1\) for divisors outside the deleted locus. Every \(m_D\)
is finite. Indeed, if a component \(E\) of \(c^*D\) dominates
\(D\) with multiplicity \(e>0\), a small transverse disc at a
general point of \(E\) has boundary mapping to a conjugate of the
\(e\)-th power of the meridian. This boundary contracts in the
compact source, so its image under \(\rho\) is trivial.

\begin{definition}\label{de:inertial-boundary}
The \emph{inertial boundary} associated with \(\rho\) on this model
of \(Z\) is
\[
 \Delta_Z=\sum_D(1-m_D^{-1})D.
\]
A \emph{prepared model} is a smooth compact K\"ahler base model,
together with a resolved source, on which the divisorial deleted locus
has simple normal crossings. If a map to a fixed torus is present,
the model is required to map holomorphically to that torus.
\end{definition}

There are only finitely many nonzero coefficients on a fixed model.
The pair \((Z,\Delta_Z)\) on a prepared model is klt. Preparation
uses resolution of the graph and discriminant; further smooth
modifications can again be prepared. The orders on all these models
refer to the same quotient \(R\). In particular, the strict transform
of a divisor has the same order. No invariance of
\(\kappa(Z,K_Z+\Delta_Z)\) under modification is assumed.
A collection of prepared models is \emph{cofinal} if every prepared
model is dominated by a member through a holomorphic modification.

\begin{lemma}[Inertia over a smooth intermediate map]
\label{de:inertia-divisibility}
Let \(f:X\to Z\) be a dominant holomorphic map of connected smooth
complex manifolds. Suppose that a homomorphism
\(\rho:\pi_1(X)\to R\) is compatible, over a dense analytic
Zariski open \(Z^\circ\), with a homomorphism
\(\psi:\pi_1(Z^\circ)\to R\). Let \(p:Z\to W\) be a smooth
proper surjective map with connected fibres, with \(W\) smooth.
Assume that the nontrivial divisorial meridian orders of \(\psi\)
give the boundary
\[
 \Delta_Z=p^*\Delta_W,
 \qquad
 \Delta_W=\sum_D(1-m_D^{-1})D,
 \qquad m_D\in\Z_{\ge2}.
\]
Thus every other divisorial meridian has trivial image. If a prime
divisor \(E\subset X\) dominates \(D\subset W\) under
\(p\circ f\), then
\begin{equation}\label{de:all-component-divisibility}
 m_D\mid \operatorname{ord}_E\bigl((p\circ f)^*D\bigr).
\end{equation}
This includes the case in which \(f(E)\) is a proper subvariety of
\(p^{-1}(D)\).
\end{lemma}

\begin{proof}
Take a general point \(w\in D\), avoiding its singularities and the
other components of \(\Delta_W\). The divisor \(p^{-1}(D)\) is
smooth near every point over \(w\). Fix such a point \(z\), and a
small coordinate ball \(U\) around it. In this ball the only
nontrivial inertial support is the smooth hypersurface
\(p^{-1}(D)\cap U\).

The map from the fundamental group of the original deleted-locus
complement in \(U\) to
\(\pi_1(U\setminus p^{-1}(D))\) is surjective. Its kernel is
normally generated by meridians around the other deleted divisors:
a transverse perturbation of a null homotopy gives precisely these
meridians. Their \(R\)-images are trivial. Removing the remaining
analytic subsets of complex codimension at least two does not change
this fundamental group, by the same transverse perturbation applied
to loops and homotopies. Consequently \(\psi\) factors locally
through
\[
 \pi_1(U\setminus p^{-1}(D))\cong\Z.
\]

At a general point of \(E\) over \(w\), choose a transverse disc
whose punctured part avoids all deleted inverse images. Its image
winds
\(e=\operatorname{ord}_E((p\circ f)^*D)\) times around the one
remaining hypersurface. Its boundary contracts in \(X\), so
compatibility with \(\rho\) gives \(\psi(\mu)^e=1\). The
order of \(\psi(\mu)\) is \(m_D\), proving
\eqref{de:all-component-divisibility}.
\end{proof}

The next consequence will also be used directly for the Albanese map.
It handles the source divisors whose images have codimension at least
two by decreasing the boundary on a fixed torus.

\begin{lemma}[A small ample boundary on a torus]
\label{de:small-boundary}
Let \(u:X\to B\) be a surjective holomorphic map from a smooth
compact K\"ahler manifold to a positive-dimensional complex torus.
Let
\(\Delta=\sum_D(1-m_D^{-1})D\) be an ample rational divisor on
\(B\), with integers \(m_D\ge2\). Suppose that
\[
 m_D\mid\operatorname{ord}_E(u^*D)
\]
for every prime divisor \(E\) dominating \(D\) under \(u\).
Then \(X\) is not special.
\end{lemma}

\begin{proof}
Put \(r=\dim B\), and let \(\cL\subset\Omega_X^r\) be the
saturated differential line of \(u\). For a local nonvanishing
generator \(\eta\) of \(K_B\), write \(a_E\) for the common
divisorial vanishing order of \(u^*\eta\) along \(E\), measured
in \(\cL\), and put \(v_E=\operatorname{ord}_E(u^*\Delta)\).
If \(E\) dominates a component \(D\), with multiplicity \(e\),
the local differential of a defining function of \(D\) gives
\[
 a_E\ge e-1\ge e(1-m_D^{-1})=v_E.
\]
If \(u(E)\) has codimension at least two, the tangent rank along
\(E\) is at most \(r-2\), and the normal direction increases
it by at most one. Thus every maximal minor of \(\dd u\)
vanishes along \(E\), and \(a_E\ge1\).

The pullback of the fixed divisor \(\Delta\) has finitely many
prime components on \(X\). We may therefore choose a rational
\(\epsilon>0\) such that \(\epsilon\le1\) and
\[
 \epsilon v_E\le a_E
 \quad\text{for every prime divisor }E\subset X.
\]
Only components with \(v_E>0\) impose a condition. This choice also
handles components whose images lie in intersections of boundary
divisors, since their \(v_E\) includes the sum of all pullback
coefficients.

The torus has trivial canonical bundle. For divisible \(m\), pullback
of a section of \(m\epsilon\Delta\), multiplied by the \(m\)-th
power of its top differential, is thus a holomorphic section of
\(\cL^{\otimes m}\). The inequality above checks every divisorial
pole; extension across codimension two finishes the check. The ratios
of these sections retain the \(r\) parameters of the ample linear
systems on \(B\). Hence \(\kappa(X,\cL)\ge r\), and
Theorem~\ref{thm:specialness-inputs} gives equality. This contradicts
specialness.
\end{proof}

\begin{remark}\label{de:fixed-map-vanishing}
The vanishing just used concerns the top differential of the fixed map
\(u:X\to B\). Saturation records its common vanishing in the map
\(u^*K_B\to\cL\). It does not assert that arbitrary lower-degree
forms vanish on exceptional divisors of a modification; that distinction
is relevant to \cite[Remark~2.2]{CampanaCorrection}.
\end{remark}

We will use the corresponding pullback statement on a neat model. A
prime divisor is \(u\)-exceptional when its image has codimension at
least two. The needed preparation is a smooth fibration model
\(u:\widehat X\to W\), with a modification
\(\nu:\widehat X\to X\) to a smooth source model, such that every
\(u\)-exceptional divisor is \(\nu\)-exceptional. Such models
exist after modifying the base; see
\cite[Lemma~1.4]{JWang2021}, whose base and source modifications
are projective morphisms. A projective initial base therefore remains
projective.
Concretely, flatten the strict transform over a modification of the
base and then resolve it. The flat strict transform has no divisors
mapping into codimension at least two. Any such divisor introduced by
the final resolution is exceptional over the original source.

\begin{lemma}[Pullback on a neat model]\label{de:neat-pullback}
Let \(u:\widehat X\to W\) and \(\nu:\widehat X\to X\)
be as above, with \(X\) smooth compact K\"ahler and \(W\) smooth
of dimension \(r>0\). Let
\(\Delta_W=\sum_D(1-m_D^{-1})D\), with \(m_D\ge2\).
Suppose that \(K_W+\Delta_W\) is big and that \(m_D\) divides
\(\operatorname{ord}_E(u^*D)\) for every component \(E\)
dominating \(D\). Then \(X\) is not special.
\end{lemma}

\begin{proof}
Let \(\cL\subset\Omega_X^r\) be the saturated differential
line of the meromorphic map induced by \(u\). A section of a
divisible multiple of \(K_W+\Delta_W\) pulls back as a meromorphic
section of \(\cL^{\otimes m}\). Along a component of multiplicity
\(e\) over \(D\), its allowed pole is cancelled because
\begin{equation}\label{de:differential-cancellation}
 e-1\ge e(1-m_D^{-1}).
\end{equation}
Divisors dominating \(W\) produce no boundary pole. The remaining
source divisors are \(u\)-exceptional and hence disappear in
codimension at least two on \(X\). Every prime divisor of \(X\)
therefore satisfies the holomorphicity condition. Extension across
codimension two gives
\[
 H^0\bigl(W,m(K_W+\Delta_W)\bigr)
 \hookrightarrow H^0\bigl(X,\cL^{\otimes m}\bigr).
\]
Pullback preserves the ratios of sections because the map is dominant.
Bigness gives \(\kappa(X,\cL)\ge r\), and
Theorem~\ref{thm:specialness-inputs} again gives the forbidden equality.
\end{proof}

\subsection{Two comparisons of models}

\begin{lemma}\label{de:birational-pushdown}
Let \(\mu:Z'\to Z\) be a modification between smooth base models
of the same fibration and quotient. For sufficiently divisible \(m\),
there is a natural injection
\[
 H^0\bigl(Z',m(K_{Z'}+\Delta_{Z'})\bigr)
 \hookrightarrow H^0\bigl(Z,m(K_Z+\Delta_Z)\bigr).
\]
\end{lemma}

\begin{proof}
View a section upstairs as a meromorphic pluricanonical form in the
common function field. Every prime divisor of \(Z\) has a strict
transform with the same meridian order. Its divisorial pole bound is
therefore the required bound on \(Z\). Extension across codimension
two gives the section downstairs, and the function-field identification
makes this map injective.
\end{proof}

Now suppose \(c:M\dashrightarrow Z\) lies over a torus \(A\),
and let \(B_c\) be the image of its general fibre group. Use
\(R=\pi_1(M)/B_c\). An isogeny \(A'\to A\) induces a subgroup
\(G'\le G=\pi_1(M)\) containing the kernel of
\(G\to\pi_1(A)\), hence containing \(B_c\). General
\(c\)-fibres lift unchanged to the source cover. The induced
quotient is consequently
\begin{equation}\label{de:quotient-inclusion}
 R'=G'/B_c\ \hookrightarrow\ G/B_c=R.
\end{equation}
On corresponding pulled-back base models, meridians lift with winding
number one. Their lattice image is zero, so their images lie in
\(R'\), and \eqref{de:quotient-inclusion} preserves their orders.
Thus both the canonical divisor and the true inertial boundary pull
back under this finite unramified base change.

\begin{lemma}[Restriction to common general fibres]
\label{de:common-fibres}
Let \(g:Z\to A\) and \(h:Z\to W\) be surjective holomorphic
maps with connected general fibres between compact complex manifolds,
with \(Z\) smooth. Let \(L\)
be a rational line bundle whose restriction to a very general smooth
\(g\)-fibre is big. Then \(L\) restricts big to every component
of a general fibre of the map
\[
 (h,g):Z\longrightarrow J\subseteq W\times A,
\]
where \(J\) denotes its image. Equivalently, one may first take a
very general \(h\)-fibre \(H\), and then a general fibre of
\(g|_H\) over its image. No equality \(g(H)=A\) is required.
\end{lemma}

\begin{proof}
For each divisible \(m\), properness gives the coherent direct image
\(g_*\cO_Z(mL)\). Choose a very general smooth fibre where base
change holds for all these countably many sheaves. Bigness on that
fibre supplies one \(m\) for which its section map has full generic
rank. On the open set where the direct image is locally free and
base change holds, evaluation defines the relative meromorphic map
to its projective bundle. The locus where its vertical differential
has full rank contains a dense analytic Zariski open \(U\subset Z\).
This condition is intrinsic to the relative evaluation and jet maps,
so it is defined over the base open, independently of a local frame
of the direct image.

A general fibre component of \(Z\to J\) meets \(U\). Otherwise
the proper analytic subset \(Z\setminus U\) would contain a
general component of a family of fibres covering \(Z\), contrary
to the fibre-dimension theorem. On such a component the relative
section map still has injective differential at a general point:
its tangent space is contained in the vertical tangent space of
\(g\). The restrictions of these sections therefore define a
generically finite meromorphic map, proving bigness.
General choices in \(J\), followed by general choices in its
projection to \(W\), give the last formulation. All assertions can
be restricted to the smooth loci supplied by generic smoothness.
\end{proof}

\Needspace{22\baselineskip}
\begin{samepage}
\subsection{Descent of positivity}

\begin{theorem}[Descent over a torus]\label{thm:torus-descent}
Let \(M\) be a special smooth compact K\"ahler manifold, and let
\(\alpha:M\to A\) be a surjective map to a complex torus with
connected fibres. Suppose that
\[
 M\dashrightarrow Z\overset{g}{\longrightarrow}A
\]
is a factorization by fibrations, and that
\(d=\dim Z-\dim A>0\). On every prepared model of \(Z\),
take the inertial boundary for
\[
 R=\pi_1(M)/\im\bigl(\pi_1(F_c)\to\pi_1(M)\bigr),
\]
where \(F_c\) is a smooth general fibre of \(M\dashrightarrow Z\).
It is impossible that, on a cofinal collection of prepared models,
\(K_Y+\Delta_Z|_Y\) is big for a very general \(g\)-fibre \(Y\).
The torus \(A\) is allowed to be a point.
\end{theorem}
\end{samepage}

\begin{proof}
Assume the stated bigness on the cofinal collection. We may further
prepare a model whenever necessary. By
Theorem~\ref{thm:wang-torus}, every model in this collection satisfies
\begin{equation}\label{de:minimum}
 \kappa(Z,K_Z+\Delta_Z)\ge d.
\end{equation}
Choose a model \(Z_0\) for which this integer is minimal, and put
\(b=\kappa(Z_0,K_{Z_0}+\Delta_{Z_0})\). Then \(b\ge d>0\).

\medskip\noindent
\emph{The Iitaka fibres.}
Resolve the Iitaka map on a modification \(\mu:Z_1\to Z_0\),
obtaining a fibration \(h:Z_1\to W\) with \(W\) smooth
projective of dimension \(b\). For this step introduce an auxiliary
klt boundary \(\Delta_1\). It has the strict-transform coefficients
of \(\Delta_{Z_0}\) and exceptional coefficients sufficiently close
to one to dominate the true inertial coefficients on \(Z_1\).
They may also be chosen so that
\begin{equation}\label{de:auxiliary-boundary}
 K_{Z_1}+\Delta_1
 =\mu^*(K_{Z_0}+\Delta_{Z_0})+E,
 \qquad E\ge0\text{ is }\mu\text{-exceptional}.
\end{equation}
Indeed all discrepancies of the log-smooth klt pair are greater than
\(-1\); an exceptional coefficient less than one can be chosen above
the negatives of these discrepancies and above the corresponding
inertial coefficient. A simultaneous log resolution makes all the
supports simple normal crossing.

Equation~\eqref{de:auxiliary-boundary} preserves sections in divisible
degrees, hence the Iitaka dimension and map. A very general smooth
\(h\)-fibre \(H\) consequently satisfies
\begin{equation}\label{de:iitaka-fibre}
 \kappa\bigl(H,K_H+\Delta_1|_H\bigr)=0.
\end{equation}
Here we use the ordinary Iitaka-fibre property, as in
\cite[Section~6.3]{JWang2021}. It can also be seen directly in this
setting. Write \(L=K_{Z_1}+\Delta_1\). By Kodaira's lemma on the projective
Iitaka base, after taking a further multiple and absorbing the
pullback of the effective remainder, the resolved system gives \(aL\sim h^*B+E\) for some divisible
\(a>0\), with \(B\) ample on \(W\) and \(E\ge0\). If a multiple \(jL|_H\) had a
positive-dimensional section map on a very general fibre, coherent
direct image and base change, followed by a sufficiently large ample
twist \(h^*(kB)\), would extend sections whose ratios vary on that
fibre. Multiplication by the canonical section of \(kE\) turns
them into sections of \((j+ka)L\). This fibrewise variation,
together with the \(b\) base parameters supplied by \(B\),
contradicts \(\kappa(Z_1,L)=b\). The restriction has a nonzero
section because a general fibre is not contained in \(E\), so its
Iitaka dimension is zero. Adjunction identifies \(L|_H\) with the
divisor appearing in \eqref{de:iitaka-fibre}.

We can arrange that a very general fibre of \(g:Z_1\to A\) has big
restriction of \(K_{Z_1}+\Delta_1\). For example, dominate
\(Z_1\) by a prepared member of the cofinal collection, perform
the same auxiliary-boundary construction relative to \(Z_0\), and
replace \(Z_1\) by this model. Its map to \(W\) is already
holomorphic. The auxiliary boundary dominates its true boundary,
whose fibre adjoint is big by assumption. Any further log resolution
preserves this fibre bigness by pullback and an effective exceptional
remainder, and \eqref{de:auxiliary-boundary} still holds.
Lemma~\ref{de:common-fibres} now makes the log canonical divisor big
on general components of common fibres of \(h\) and \(g\).
Generic smoothness, also for the finitely many boundary strata,
justifies the adjunction restrictions to these common fibres.

Apply Theorem~\ref{thm:wang-torus} to the klt pair in
\eqref{de:iitaka-fibre}. Its Albanese map is surjective with connected
fibres. The map \(g|_H\) factors through a homomorphism from
\(\Alb(H)\) to \(A\), followed by a translation. Quotienting by
the connected kernel shows that its Stein base is a torus finite
\'etale over its image subtorus translate. Torus subadditivity
on \(H\), using the bigness of the common fibres, forces their
dimension to be zero. The Albanese map of \(H\) is therefore
birational, and \(g|_H\) is, birationally, an isogeny onto a
subtorus translate.

Every component of \(\Delta_1|_H\) is exceptional for this
birational Albanese map. To prove this, suppose that a component maps
onto a divisor \(D\) of the torus. The relative canonical divisor
of the birational map to the smooth torus is effective, with positive
coefficients on all exceptional primes. For sufficiently small
rational \(\epsilon>0\), it absorbs the exceptional components of
the pullback of \(\epsilon D\), while the chosen boundary
component supplies its strict transform. Thus
\(K_H+\Delta_1|_H\) contains that pullback up to rational linear
equivalence. By Lemma~\ref{lem:torus-divisor}, a nonzero effective
divisor on a complex torus has positive Iitaka dimension. This contradicts
\eqref{de:iitaka-fibre}.

\medskip\noindent
\emph{A torus-bundle model.}
On the connected smooth parameter open of \(h\), the image of
\(H_1(H,\Z)\) in \(H_1(A,\Z)\) is locally constant. Transport
in the proper smooth family and the total map to \(A\) identify
these images. Their real span determines a fixed subtorus
\(T\subset A\), so the image of a general \(H\) is a translate
of \(T\). Write \(\Lambda_T\) for its lattice and \(L_H\)
for the finite-index image lattice. Choose \(k>0\) with
\(k\Lambda_T\subseteq L_H\), and use the isogeny
\(A_1\to A\) given by multiplication by \(k\).

Let \(T_1\) be the identity component of the inverse image of
\(T\). The lattice of a subtorus is primitive in the ambient
lattice, so the image of the lattice of \(T_1\) is
\(k\Lambda_T\). The choice of \(k\) makes each connected lifted
general \(H\)-fibre map with degree one onto a translate of
\(T_1\). This is a lattice construction; no splitting of the torus
extension is required.

Pull back both the source and \(Z_1\) to \(A_1\). These covers
are connected, since the original maps to \(A\) are surjective
with connected fibres and therefore surject on \(\pi_1(A)\).
Write \(M_1\) for the source cover. The Stein factorization of
the pulled-back map to \(W\) has a finite base over \(W\); take
a smooth projective model \(W_1\) of that base. The map to
\(A_1/T_1\) is constant on its general fibres and hence descends
meromorphically to \(W_1\). Resolve this map as well. The pair of
projection maps then gives a birational model
\begin{equation}\label{de:torus-bundle}
 Z_*=W_1\times_{A_1/T_1}A_1,
 \qquad p:Z_*\longrightarrow W_1.
\end{equation}
Indeed it has degree one on the general fibres just described.
The space \(Z_*\) is smooth, and it is a closed submanifold of the
compact K\"ahler product \(W_1\times A_1\). The map \(p\) is
a principal holomorphic bundle with connected torus fibre \(T_1\).
Its transition functions are translations, so an invariant top form
on \(T_1\) trivializes the relative canonical bundle. In particular,
\begin{equation}\label{de:canonical-pullback}
 K_{Z_*}=p^*K_{W_1}.
\end{equation}
We are free to replace \(W_1\) by any further smooth projective modification
and use its corresponding fibre product in
\eqref{de:torus-bundle}.

The true inertia on \(Z_*\) has no horizontal divisor over
\(W_1\). To check this on a morphism of models, resolve the maps
from the pulled-back \(Z_1\) to \(W_1\) and \(Z_*\), and
carry the auxiliary boundary by the rule
\eqref{de:auxiliary-boundary}. Its restriction to a general fibre
still has only Albanese-exceptional components. Indeed finite
\'etale covers of a birational torus are birational to torus
covers. An exceptional prime on a further model either misses the
general fibre or restricts to a divisor exceptional over the old
general fibre, so it remains Albanese-exceptional. The strict transform
of a horizontal prime of
\(Z_*\) maps onto a divisor of its general torus fibre. It cannot
occur in this auxiliary boundary, which dominates true inertia.
Its meridian order is therefore one.

Every vertical prime divisor of the smooth bundle \(p\) is the
full inverse image of a prime divisor of \(W_1\). A divisor cannot
map into codimension at least two by the fibre-dimension formula,
and connected smooth fibres are irreducible. We have consequently
obtained a rational boundary \(\Delta_{W_1}\) with
\begin{equation}\label{de:boundary-pullback}
 \Delta_{Z_*}=p^*\Delta_{W_1},
\end{equation}
with the same finite integral meridian orders on corresponding
divisors.

\medskip\noindent
\emph{The true boundary is big downstairs.}
Let \(F\) denote the finite deck group of \(A_1\to A\).
Write \(\widetilde Z_1=Z_1\times_A A_1\) and let
\(\phi:\widetilde Z_1\dashrightarrow Z_*\) be the birational
comparison. Take the closure of the simultaneous graph of
\(\phi\circ f\), for \(f\in F\), in
\(\widetilde Z_1\times (Z_*)^F\). The group acts on this graph
by its action on \(\widetilde Z_1\) and permutation of the other
factors. Resolve equivariantly and take a further equivariant
K\"ahler modification using
\cite[Theorem~1.1 in the author's version]{JiaMeng2024}; the
finite group preserves the average of a big class. This gives a
smooth K\"ahler \(V\) mapping to \(\widetilde Z_1\), to
\(Z_*\), and to \(A_1\).
The action is free because its map to \(A_1\) is equivariant for
the free translation action. The quotient \(V/F\) is thus smooth
K\"ahler and is a modification of \(Z_1\), hence of \(Z_0\),
over the original torus \(A\). Moreover, \(V\to V/F\) is the
pullback of \(A_1\to A\): each deck orbit maps bijectively onto
the corresponding isogeny fibre.

Further prepare this quotient, and, if necessary, dominate it by
a member of the cofinal collection. Pull this modification back
along \(A_1\to A\); all comparison morphisms persist. We continue
to call the resulting cover \(V\). Its true inertial adjoint is
the unramified pullback of that on its downstairs quotient, by
\eqref{de:quotient-inclusion}. The minimum in
\eqref{de:minimum} therefore gives
\[
 \kappa(V,K_V+\Delta_V)\ge b.
\]
All boundaries in this inequality are true inertial boundaries;
the auxiliary exceptional coefficients of \(\Delta_1\) are not
used here. Lemma~\ref{de:birational-pushdown} applied to
\(V\to Z_*\) gives the same lower bound on \(Z_*\).
Equations~\eqref{de:canonical-pullback} and
\eqref{de:boundary-pullback}, together with \(p_*\cO_{Z_*}=
\cO_{W_1}\), now yield
\[
 \kappa(W_1,K_{W_1}+\Delta_{W_1})\ge b=\dim W_1.
\]
Thus \(K_{W_1}+\Delta_{W_1}\) is big.

This argument applies after any chosen further smooth projective modification
of \(W_1\): form its fibre product, repeat the equivariant
comparison, and use the same minimum on the prepared quotient over
\(A\). The newly appearing true coefficients are read on the new
model. No comparison of exceptional coefficients on unrelated models
is needed.

\medskip\noindent
\emph{The differential contradiction.}
Choose \(W_1\) sufficiently modified that the induced fibration
from a resolution of \(M_1\) to \(W_1\) is neat, as described
before Lemma~\ref{de:neat-pullback}. The preceding argument still
gives bigness on this choice. The source resolution maps holomorphically
to \(Z_*\), using its maps to \(W_1\) and \(A_1\). Its
fundamental-group quotient is the one used to define the true inertia
on \(Z_*\). Lemma~\ref{de:inertia-divisibility}, with the smooth
map \(p\), shows that every boundary order on \(W_1\) divides
the multiplicity of every source component dominating that divisor.
Lemma~\ref{de:neat-pullback} therefore gives a Bogomolov sheaf on a
smooth source model of \(M_1\). This contradicts
Theorem~\ref{thm:specialness-inputs}, since \(M_1\) is a connected
finite \'etale cover of the special manifold \(M\).
\end{proof}
